\section{Experiments}

\subsection{Setup}

All experiments are built on SemlaFlow, one of the three mixed-domain flow matching models
described in Section~\ref{sec:existing-models}. It was chosen because it is fully open source,
uses the endpoint prediction parameterisation of Section~\ref{sec:flow-matching} (so the
averaged target of Section~\ref{sec:research-gap} is exactly the quantity its coordinate loss
already regresses towards), and because it isolates Hungarian alignment as a single step in its
training loop, meaning permutation averaging can be introduced by replacing one function.

\paragraph{The four training variants.}
We compare four ways of choosing the pairing between noise and data, and nothing else. Given a data
molecule $z$, a noise sample $x$ and a time $t$, the variants are:

\begin{itemize}
    \item \textbf{None.} No alignment. The interpolant is built directly from the sampled pair and
        the target is $z$ itself. This is the independent-coupling floor.
    \item \textbf{Hungarian.} Before interpolating, the noise is permuted onto the data: the
        Hungarian algorithm is run on $C_{ij} = \|z_i - x_j\|^2$ to give $\sigma^*$, and the
        interpolant is built from $P_{\sigma^*} x$ rather than $x$. The target remains $z$. This is
        the standard practice described in Section~\ref{sec:hungarian-kabsch} and is the published
        baseline.
    \item \textbf{Sinkhorn.} The interpolant is built from the sampled pair as in \textsc{None},
        and the \emph{target} is then replaced by $\bar{P} z$, where $\bar{P}$ is the doubly
        stochastic matrix returned by the Sinkhorn algorithm on the cost matrix and temperature of
        Equation~\eqref{eq:sinkhorn-temperature}. This is the deterministic approximation to
        permutation averaging of Section~\ref{sec:sinkhorn-approx}.
    \item \textbf{Metropolis.} As \textsc{Sinkhorn}, but the target is $P_\sigma z$ for a single
        permutation $\sigma$ drawn from $q_t(\sigma \mid x_t, z)$ by the Markov chain of
        Section~\ref{sec:metropolis-sampling}. This is the unbiased stochastic counterpart to
        Sinkhorn.
\end{itemize}

Two differences between these are worth making explicit, because they are easy to miss.

First, \emph{which side} the permutation is applied to. Hungarian alignment as implemented permutes
the noise, but it could equally be written as permuting the data, and the two give the same
objective. Since $P_{\sigma^*}$ is orthogonal,
$\|P_{\sigma^*} x - z\|^2 = \|x - P_{\sigma^*}^\top z\|^2$, so the two minimisation problems
correspond. Relabelling every atom slot of the first by $P_{\sigma^*}^\top$ gives
$P_{\sigma^*}^\top x_t = (1-t)x + t P_{\sigma^*}^\top z$ with target $P_{\sigma^*}^\top z$, which is
the second; and since $f^\theta$ is permutation-equivariant and the loss is a sum over atoms (and,
for bonds, over pairs, with $e \mapsto P e P^\top$), relabelling both the prediction and the target
by the same permutation leaves the loss unchanged. Hungarian may therefore be read either way.

This is a statement about which \emph{side} the permutation acts on, and not about \emph{when} it is
computed --- those are different questions, and the answer to the second is given below.

Second, \emph{what} the permutation is computed from, and here the choice genuinely matters. The
Hungarian cost compares the raw noise with the data, since at that point no interpolant exists yet.
The Sinkhorn and Metropolis costs compare $x_t$ with $tz$, as required by the posterior of
Equation~\eqref{eq:permutation-posterior}, which is only defined given a point part-way along the
flow. These are not the same assignment problem. Discarding the terms that do not depend on
$\sigma$, minimising either cost is equivalent to maximising an inner-product score, and
substituting $x_t = (1-t)x + tz$ relates the two exactly:
\begin{equation}
    S_{x_t}(\sigma) = (1-t)\, S_{x}(\sigma) \;+\; t\, S_{z}(\sigma),
    \qquad
    S_{a}(\sigma) := \sum_{i=1}^{N} \langle a_i,\, z_{\sigma(i)} \rangle .
    \label{eq:cost-decomposition}
\end{equation}
The second term is maximised at the identity: by Cauchy--Schwarz
$\langle z_i, z_{\sigma(i)}\rangle \le \|z_i\|\|z_{\sigma(i)}\|$, and
$\sum_i \|z_i\|\|z_{\sigma(i)}\| \le \sum_i \|z_i\|^2$ for any permutation, again by
Cauchy--Schwarz. So the cost built from $x_t$ is the cost built from $x$ plus a bias towards the
identity of relative weight $t/(1-t)$. The two coincide only in the limit $t \to 0$, where the cost
becomes uniform and every permutation is equally good, consistent with
Section~\ref{sec:sinkhorn-approx}. The consequence of this is taken up in the discussion.

A consequence of the first point is that all four models are trained to transport the same prior. A
permutation of independent Gaussian noise is still independent Gaussian noise, so the Hungarian arm
does not change the distribution its network learns to start from; and the averaging arms do not
touch the noise at all. The four models are therefore directly comparable at generation time, where
each is simulated from an ordinary Gaussian sample. Note that this would \emph{not} hold if
$\bar{P}$ were applied to the noise rather than the target: a doubly stochastic matrix is a convex
combination of permutations rather than a permutation, so averaging the noise would shrink its
variance and train the model on a prior that does not exist at sampling time.

\begin{algorithm}
\caption{Permutation-Averaged Flow Matching Training}
\label{alg:permutation-averaged-training}
\begin{algorithmic}[1]
\Require Dataset $z \sim p_\mathcal{D}$, priors $p_{\text{init}}$ for each variable, network
    $f^\theta$, method $\in \{\textsc{None}, \textsc{Hungarian}, \textsc{Sinkhorn},
    \textsc{Metropolis}\}$.
\For{each mini-batch of data}
    \State Sample a molecule $z = (r, a, c, e) \sim p_\mathcal{D}$ with $N$ atoms
    \State Sample a time $t \sim \mathcal{U}(0,1)$
    \State Sample an initial molecule $x = (r_0, a_0, c_0, e_0)$, each component from its prior
    \If{method $=$ \textsc{Hungarian}} \Comment{align the noise, before interpolating}
        \State $\sigma^* \gets$ Hungarian solution of $C_{ij} = \|z_i - x_j\|^2$
        \State $x \gets P_{\sigma^*} x$
    \EndIf
    \State Rotate $x$ onto $z$ with the Kabsch algorithm
    \State Set $x_t \gets (1-t)\,x + t\,z$, blending coordinates and resampling each categorical
        variable from $z$ with probability $\kappa_t$
    \State $z' \gets z$ \Comment{the target, modified only by the averaging methods}
    \If{method $\in \{\textsc{Sinkhorn}, \textsc{Metropolis}\}$}
        \State $C_{ij} \gets \|(x_t)_i - t z_j\|^2$, \quad
            $\varepsilon_t \gets 2(1-t)^2$
        \If{method $=$ \textsc{Sinkhorn}}
            \State $\bar{P} \gets \textsc{Sinkhorn}(C, \varepsilon_t)$
        \Else
            \State $\bar{P} \gets P_\sigma$, $\sigma$ a sample from $q_t(\cdot \mid x_t, z)$
        \EndIf
        \State $z' \gets (\bar{P}r,\; \bar{P}a,\; \bar{P}c,\; \bar{P}e\bar{P}^\top)$
    \EndIf
    \State Compute the loss of Equation~\eqref{eq:mixed-domain-loss} using $z'$ in place of $z$
    \State Update parameters $\theta \gets \mathrm{grad\_update}(\mathcal{L}(\theta))$
\EndFor
\end{algorithmic}
\end{algorithm}

\paragraph{Method parameters.}
The Sinkhorn iteration is run for $100$ steps, which is enough for the row and column sums to
agree with their targets to within $10^{-2}$ at every $t$. The Metropolis chain is run for $100$
steps. Two details differ slightly from the description in Section~\ref{sec:metropolis-sampling}.
First, the chain is initialised at the identity permutation rather than at the Hungarian solution:
because $x_t$ is constructed from $z$ in its given atom order, the identity \emph{is} the
permutation that produced $x_t$, so it is already a high-probability starting point and is
available without solving an assignment problem inside the training loop. Second, rather than
imposing a distance threshold $d_{\max}$, proposals are restricted to swaps between each atom and
its $8$ nearest neighbours in $x_t$; this has the same effect of suppressing proposals that would
be rejected almost surely, but does not require a threshold to be tuned per dataset.

The temperature $\varepsilon_t$ carries no free hyperparameter, being fixed by the variance of the
conditional path. This is deliberate: it means the Sinkhorn arm cannot be made to look better or
worse by tuning, and any difference from the Hungarian arm is attributable to the averaging itself.

\paragraph{What is held fixed.}
Kabsch alignment is enabled in all four arms. This is the harder test, since any benefit from
permutation averaging must appear on top of rotation alignment rather than as a substitute for it,
and it keeps the Hungarian arm faithful to the published baseline. We do \emph{not} experiment with
rotation averaging in the de novo setting, even though the results of
Section~\ref{sec:rotation-averaging} have not previously been transferred to this domain; and we do
not attempt to combine rotation and permutation averaging, which is not straightforward because the
optimal rotation and the optimal permutation depend on one another. Both are left for future work
contingent on permutation averaging succeeding in isolation. Architecture, optimiser, data order,
noise draws and sampling schedule are identical across arms, and all four use the same random seed.

\paragraph{Training.}
All models are trained on QM9 for $300$ epochs at a constant learning rate of $3\times 10^{-4}$
after $2000$ warm-up steps, with gradients clipped at norm $1$. Times are sampled uniformly,
$t \sim \mathrm{Beta}(1,1)$. The loss weights of Equation~\eqref{eq:mixed-domain-loss} are
$\lambda_a = 0.2$ for atom types, $\lambda_c = 1.0$ for charges and $\lambda_e = 0.5$ for bonds,
with the coordinate term at unit weight. These are the reference values for this architecture,
of which only the bond weight is specific to QM9 --- it is halved from the value used for
GEOM-Drugs, since QM9's molecules are small enough that the bond term would otherwise dominate.
Each run takes approximately six to seven hours on a single GH200 GPU.

\paragraph{Generation.}
Molecules are generated by simulating the joint process of Algorithm~\ref{alg:joint-simulation}
with the Euler method on a uniform grid of $100$ steps over $[0,1]$, so $h = 0.01$. The coordinate
update is $r_{t+h} = r_t + h\,v^\theta_t$, with the vector field recovered from the endpoint
prediction by Equation~\eqref{eq:endpoint-to-vector-field}, which for $\alpha_t = t$,
$\beta_t = 1-t$ is simply $v^\theta_t = (r^\theta_t(w_t) - r_t)/(1-t)$. The atom type, charge and
bond components are updated by the corresponding jump step of Equation~\eqref{eq:dfm-jump} on the
same grid. Simulation is otherwise deterministic: no noise is added to the coordinate update. Since
the network is called once per step, $100$ steps corresponds to an NFE of $100$.

\paragraph{Evaluation.}
Each trained model generates $5000$ molecules. We report validity and
molecule stability computed with the corrected valency table discussed in
Section~\ref{sec:evaluation}, and, for the 3D domain, GFN2-xTB rather than MMFF: each generated
molecule is relaxed and we record the relaxation energy $\Delta E_{\text{relax}}$ together with the
mean bond length, bond angle and torsion deviations between the generated molecule and its own
relaxed counterpart, plus the RMSD between the two. Following the recommendation in
Section~\ref{sec:evaluation}, each of these is reported as both a median and a mean, since the
distributions are right-skewed and the two capture different failure modes: the median describes
typical geometry quality, the mean the rate of catastrophic failures. Uniqueness and novelty are
not reported, for the reasons given there.

We also report a diagnostic requiring no trained model at all: the ratio
$\|\bar{P}z\| / \|z\|$ measured on real QM9 molecules as a function of $t$, together with the
effective number of atoms $1 / \sum_j \bar{P}_{ij}^2$ being averaged into each row of the target.
These describe the training target itself and are used in the discussion.

\subsection{Results}

Table~\ref{tab:qm9-results} gives the headline numbers. We report QM9 only. The failure described
below is severe enough that training on GEOM-Drugs would not be informative: a method that destroys
the smaller and easier dataset is extremely unlikely to succeed on the larger and harder one.

\begin{table}[h]
\centering
\small
\begin{tabular}{lcccccc}
\hline
 & & $\Delta E_{\text{relax}}$ & Bond length & Bond angle & Torsion & RMSD \\
Method & Validity & (kcal/mol) & (\AA) & ($^\circ$) & ($^\circ$) & (\AA) \\
\hline
None        & 0.994 & 15.2 & 0.029 & 2.04 & 6.66 & 0.197 \\
Hungarian   & 0.992 & \textbf{13.5} & \textbf{0.027} & \textbf{1.88} & \textbf{5.94}
            & \textbf{0.180} \\
Sinkhorn    & 0.000 & --- & --- & --- & --- & --- \\
Metropolis  & 0.967 & 53.3 & 0.034 & 2.25 & 7.93 & 0.235 \\
\hline
Sinkhorn-HC & 0.950 & 2803 & 0.297 & 16.7 & 35.5 & 0.859 \\
\hline
\end{tabular}
\caption{QM9 results, $5000$ molecules per method at $100$ integration steps, one seed. Every 3D
column is a mean deviation from the molecule's own GFN2-xTB-relaxed counterpart, so lower is better
throughout. The relaxation succeeded and converged for between $97\%$ and $99.9\%$ of the molecules
in each set. Sinkhorn-HC is the ablation described below, in which coordinates are averaged but the
categorical targets are left as hard one-hot labels.}
\label{tab:qm9-results}
\end{table}

\paragraph{A note on the metrics reported.}
Two remarks on what Table~\ref{tab:qm9-results} does and does not contain.

Atom and molecule stability are both omitted because on QM9 they are degenerate. Atom stability is
$1.000$ for every method: under the corrected valency table, every molecule that RDKit could
construct had all of its atoms in a valid valency state. Molecule stability follows from this and
is consequently equal to validity to four decimal places in all five cases, so it would duplicate a
column already present. Validity is therefore the only 2D metric that carries any information here,
and it does not separate the four methods either --- with the single exception of the Sinkhorn arm,
which it separates only because that arm produces nothing usable at all. All of the discrimination
in the table comes from the 3D columns, which is a concrete instance of the concern raised in
Section~\ref{sec:evaluation} about 2D metrics saturating.

The 3D columns are means. All five distributions are right-skewed, so the medians are smaller
throughout, and by varying amounts: the ratio of mean to median is around $1.2$ for bond lengths
and angles, but reaches $1.9$ for RMSD and $3.7$ for the relaxation energy of the Metropolis arm,
where a small number of badly-formed molecules dominates the average. The ranking of the methods is
the same under either statistic for all five metrics, so nothing that follows depends on the
choice; where the difference between the two is itself informative it is given explicitly in the
text.

\paragraph{Neither averaging method improves on hard alignment.}
Hungarian alignment is the best method on every 3D metric, and the improvement over no alignment at
all is consistent but small (Cliff's $\delta = -0.08$ on $\Delta E_{\text{relax}}$, with a
bootstrap confidence interval on the difference of medians of $[-1.09, -0.41]$ kcal/mol). The
Metropolis target is worse than Hungarian on every metric and worse than no alignment at all
(Cliff's $\delta = +0.31$). Its relaxation energy is the clearest case where the mean and the
median disagree: $53.3$ kcal/mol against a median of $14.3$, where the baseline sits at $15.2$ and
$10.0$. The typical Metropolis molecule is therefore only modestly worse than the baseline, and the
inflated mean comes from a tail of badly-formed molecules rather than from a uniform degradation.
The Sinkhorn target fails completely: not a single one of its $5000$ generated molecules could be
constructed by RDKit.

\paragraph{The Sinkhorn failure is in the bond graph, not the geometry.}
A zero validity score reports only that the molecular graph is unusable, so it is worth asking
whether the coordinates collapsed as well. They did not. The Sinkhorn samples have a median radius
of gyration of $2.08$~\AA{} against $2.31$~\AA{} for the working methods, and a mean
nearest-neighbour distance of $1.13$~\AA{}, which is an ordinary bond length. The molecules are
therefore of roughly the right size and their atoms are appropriately spaced; what is wrong is the
predicted bond graph.

\paragraph{Isolating the categorical channel.}
Since permutation averaging changes both the coordinate target and the categorical targets, we ran
one further ablation, \textbf{Sinkhorn-HC}, in which the doubly stochastic matrix is applied to the
coordinates as before but the atom type, charge and bond targets are left as the hard one-hot
labels of the identity permutation. Validity is restored from $0.000$ to $0.950$, confirming that
the mixing of categorical targets is what makes the bond graph unconstructible.

The recovery is only apparent, however. Under GFN2-xTB the Sinkhorn-HC molecules are far worse than
any other arm: a median relaxation energy of $769$ kcal/mol against $10$ for the baseline, bond
lengths wrong by $0.25$~\AA{} rather than $0.025$~\AA{}, and bond angles wrong by $16^\circ$ rather
than $1.7^\circ$. The separation from the baseline is essentially total (Cliff's $\delta = 0.995$).
So the two channels fail independently and for different reasons: averaging the categorical targets
destroys the graph, and averaging the coordinate targets destroys the geometry. It is worth noting
that neither validity nor radius of gyration detected the second failure --- Sinkhorn-HC looks
healthy on both --- and it was only visible once a quantum-chemical method was used.

\paragraph{The averaged target collapses towards the origin.}
Table~\ref{tab:target-collapse} measures the training target directly, with no model involved. At
small $t$ the Sinkhorn plan is close to uniform, so each row of $\bar{P}z$ is close to the mean of
all $N$ atom positions. QM9 molecules are centred at the origin, so that mean is the origin, and
the regression target is close to zero.

\begin{table}[h]
\centering
\begin{tabular}{lcccccc}
\hline
$t$ & 0.05 & 0.2 & 0.5 & 0.7 & 0.9 & 0.95 \\
\hline
$\|\bar{P}z\| / \|z\|$ & 0.04 & 0.20 & 0.62 & 0.87 & 0.99 & 0.99 \\
Effective atoms averaged & 17.3 & 15.8 & 7.5 & 3.1 & 1.2 & 1.1 \\
\hline
\end{tabular}
\caption{The Sinkhorn target measured on real QM9 molecules. The ratio is the mean per-atom
distance of the target from the origin, divided by the same quantity for the unmodified data
molecule; the effective atom count is $1 / \sum_j \bar{P}_{ij}^2$. For the Hungarian and Metropolis
methods the ratio is exactly $1$ at every $t$, since both apply a genuine permutation.}
\label{tab:target-collapse}
\end{table}

The behaviour at the two ends of the flow is exactly what Section~\ref{sec:sinkhorn-approx}
predicts, and in that sense the implementation is working as designed: the plan is uniform at
$t \to 0$ and converges to a single permutation at $t \to 1$. The problem is what this implies for
training. Over the whole first half of every trajectory the model is being asked to predict
something close to the centroid of the molecule rather than the molecule, and since the endpoint
prediction is converted into a vector field by Equation~\eqref{eq:endpoint-to-vector-field}, the
early part of the simulated flow is driven by a network trained to point at the origin.

\subsection{Discussion}

Neither approximation to permutation averaging improved on hard Hungarian alignment, and the
deterministic one failed outright. Three explanations account for what we observe.

\paragraph{The cost matrix ignores the discrete domain.}
The cost $C_{ij} = \|(x_t)_i - t z_j\|^2$ is purely geometric, so a carbon paired with a hydrogen
is considered exactly as good as a carbon paired with a carbon at the same distance. Under hard
alignment this is tolerable, since the result is still a genuine relabelling of a real molecule and
the target remains a valid chemical structure whichever permutation is chosen. Under averaging it
is not: the plan spreads mass across atoms that happen to be nearby regardless of what they are,
and the resulting target asks the network to predict an atom that is part carbon and part hydrogen,
with bonds to match. The zero validity of the Sinkhorn arm and its recovery under Sinkhorn-HC show
that this is not a minor effect but the dominant one in the categorical channel. A cost that
included a penalty for mismatched atom types would concentrate the plan on chemically sensible
pairings and is the most obvious thing to try next.

\paragraph{Averaging over permutations of a known molecule is not the same as averaging over
endpoints.}
The justification for an averaged target in Section~\ref{sec:rotation-averaging} is that the
endpoint prediction is trying to learn $\mathbb{E}_{z \sim p(z \mid x_t)}[z]$, so replacing a single
representative by an average over representations moves the target towards what the loss is already
estimating. This argument requires the average to be taken over the correct conditioning set. What
$\bar{P}z$ actually computes is an average over permutations of the \emph{specific} molecule $z$
drawn in that mini-batch, which conditions on information the network does not have. The result
moves the target away from $z$ without moving it towards $\mathbb{E}[z \mid x_t]$; it moves it
towards that one molecule's centroid, which is what Table~\ref{tab:target-collapse} shows. The
target is being averaged over the wrong thing.

\paragraph{The posterior is peaked on the identity by construction.}
Equation~\eqref{eq:cost-decomposition} makes the self-consistency of the construction explicit. The
cost is built from $x_t$, but $x_t$ was itself built from $z$ in its given atom order, so the
identity is the permutation that produced it and the cost carries an explicit bias towards the
identity at every $t$, not merely in the limit. The two contributions have very different scales:
$S_z(\mathrm{id}) = \sum_i \|z_i\|^2$ is large and positive, whereas $S_x$ is a sum of inner
products between independent vectors and is close to zero in expectation. Measuring this on random
$18$-atom problems, the permutation minimising the $x_t$-based cost already differs from the one
minimising the $x$-based cost in more than half of cases at $t = 0.05$, and never agrees with it
beyond $t = 0.2$; by $t = 0.8$ it is the identity in $63\%$ of cases and by $t = 0.95$ in
essentially all of them.

This is not in itself an error --- $x_t$ really does carry information about which permutation
generated it, and a posterior given $x_t$ ought to reflect that. The difficulty is that the
generating permutation is always the identity here, so the posterior concentrates on it for a
reason that has nothing to do with the geometry of the molecule. Combined with the previous
paragraph, the picture is that at low $t$ the target is smeared across the whole molecule and at
high $t$ it is pinned to the permutation it started from, leaving only a narrow band of $t$ in which
the averaging is doing anything both non-trivial and informative.

\paragraph{Why the same problem does not arise for rotations.}
It is fair to ask why the rotation-averaged objective of Section~\ref{sec:rotation-averaging}
succeeds when the permutation analogue does not, since averaging a molecule over all of $SO(3)$ at
$t \to 0$ would also produce a spherically symmetric and thus uninformative target. Two differences
matter. First, $SO(3)$ is a three-dimensional compact group whereas $S_N$ has $N!$ elements, so the
permutation posterior is diffuse over a vastly larger set and concentrates far more slowly in $t$;
Table~\ref{tab:target-collapse} shows more than seventeen atoms still being averaged at $t = 0.2$.
Second, and more importantly, a rotation acts only on coordinates and leaves the discrete domain
untouched, so rotation averaging has no categorical channel to damage. Permutation averaging has
both problems and the discrete one is the more severe.

\paragraph{Is permutation averaging hopeless?}
These results do not establish that. What they establish is that this particular construction of it
fails, and they identify two specific reasons why. Both are properties of how the posterior was
defined rather than of the idea of averaging over permutations: a cost matrix blind to atom types,
and a posterior conditioned on the batch's own data molecule. Since the failure is severe rather
than marginal, however, the honest conclusion at this point is that permutation-averaged flow
matching does not work as a drop-in replacement for Hungarian alignment, and that the burden is on
a future version to fix the conditioning before the hypothesis of
Section~\ref{sec:research-gap} can be tested properly.

\paragraph{A caveat on the strength of these claims.}
Each arm was trained once, so strictly these numbers compare five trained models rather than five
methods, and run-to-run variation is not separated out. The Sinkhorn and Sinkhorn-HC effects are
far too large for this to matter. The advantage of Hungarian over no alignment, however, is small
enough that it would need further seeds before being claimed with confidence.
